%!TEX root = ../csuthesis_main.tex

{~}
\vspace{18pt}
\section{攻读学位期间主要的研究成果} % 无章节编号

\ifblindreview
% \noindent
% （盲审隐去作者相关具体信息）
\fi
\vspace{11pt}
\subsection*{一、学术论文}

\ifblindreview

% \noindent
% 第一作者：JCR 1 区 x 篇，会议 x 篇 \\{}
% 第二作者：JCR 1 区 x 篇，3 区 x 篇，4 区 x 篇，EI x 篇 

% \noindent
% 投稿状态： 
% IEEE Transactions on Image Processing 1篇（under review）\\{}
% IEEE Transactions on Circuits and Systems for Video Technology 1 篇（Accept with Minor Revision）

% 学位办老师要求用如下这种几乎算是单盲的格式，我也木有办法……
\subsubsection*{已录用/检索论文}
%\noindent
第一作者：
\begin{enumerate}[label={[\arabic*]},itemindent=2em,wide]
	\item CSU Latex Template[J]. CSU player: 1(1):1-10. {\bfseries \heiti（SCI 录用，JCR 1 区）}
	\item CSU Latex Template[J]. CSU player: 1(1):1-10. {\bfseries \heiti（SCI 检索，JCR 2 区）}
\end{enumerate}
第二作者：
\begin{enumerate}[label={[\arabic*]},itemindent=2em,wide]
	\item CSU Latex Template[J]. CSU player: 1(1):1-10. {\bfseries \heiti（SCI 检索，JCR 1 区）}
	\item CSU Latex Template[J]. CSU player: 1(1):1-10. {\bfseries \heiti（SCI 检索，JCR 2 区）}
\end{enumerate}
%\noindent
\subsubsection*{投稿状态论文}
%\noindent
第一作者：
\begin{enumerate}[label={[\arabic*]},itemindent=2em,wide]
	\item CSU Latex Template. XXX Transactions on CSU player. {\bfseries \heiti（SCI Under Review，JCR 1 区）}
\end{enumerate}
第二作者：
\begin{enumerate}[label={[\arabic*]},itemindent=2em,wide]
	\item CSU Latex Template. XXX Transactions on CSU player. {\bfseries \heiti（SCI Under Review，JCR 1 区）}
	\item CSU Latex Template. XXX Transactions on CSU player. {\bfseries \heiti（SCI Under Review，JCR 2 区）}
\end{enumerate}


\else
% 标准版本
\begin{enumerate}[label={[\arabic*]},itemindent=2em,wide]
	\item Deyu Zhang, \textbf{Kangqi Wei}, Huan Yang, Yin Tang, Fan Wu. IEEE Internet of Things Journal[J]. StereoMulti3DPose: Multi-Person 3D Pose Estimation via
Lightweight Asymmetric Stereo Learning on the Edge. {\bfseries \heiti（Minor Revision，JCR 1 区）}
\end{enumerate}
\fi

\vspace{22pt}
\subsection*{二、专利软著}
\ifblindreview
软著 1 项，已公开
\else
\begin{enumerate}[label={[\arabic*]},itemindent=2em,wide]
	\item 中南大学. 双目3D人体姿态估计系统V1.0. 登记号：2026SR0136938
\end{enumerate}
\fi

\ifblindreview
\else

% \vspace{22pt}
% \subsection*{三、主持和参与的科研项目}
% \begin{enumerate}[label={[\arabic*]},itemindent=2em,wide]
% 	\item 国家自然科学基金面上项目《XXXXXXXXXXXX》， 项目编号：XXXXXXXX，参与.
% \end{enumerate}

% {~}
\vspace{22pt}
\subsection*{三、个人获奖情况}
%\noindent
\begin{enumerate}[label={[\arabic*]},itemindent=2em,wide]
	\item 2023年中南大学一等金奖
	\item 2024年中南大学三等金奖
	\item 2025年中南大学三等金奖
\end{enumerate}
\fi

\newpage

\ifblindreview
\else

{~}
\vspace{18pt}
\section{{致~~~~谢}} % 无章节编号
% 重新设置正文行间距，因为前置部分设置时候行间距被改过
\renewcommand*{\baselinestretch}{1.0} % 几倍行间距
\setlength{\baselineskip}{20pt} % 基准行间距

论文行将收束，心中最深的感受首先是感恩。在硕士阶段的学习与科研过程中，我有幸得到许多师长、亲友与同伴的关心和帮助，而其中最需要郑重表达谢意的，是我的导师张德宇老师。

首先，我要向张德宇老师致以最诚挚、最深切的感谢。自进入课题以来，张老师始终以严谨的治学态度、开阔的学术视野和高度负责的育人精神指导我前行。从研究方向的确定、技术路线的梳理，到实验细节的打磨、论文结构的推敲，张老师都给予了我耐心而细致的指导。每当我在研究中遇到困难、在实验结果面前感到迷惘、在论文写作中陷入反复时，张老师总能以敏锐而清晰的判断帮助我抓住问题本质，引导我重新建立思路。您不仅教会了我如何做研究，更让我体会到什么是严谨求实的学术品格，什么是对学生认真负责的师者风范。

尤为令我感佩的是，张老师始终坚持高标准、严要求，但这种要求之中又饱含着耐心、包容与鼓励。您对每一个技术细节的认真、对每一处表述的推敲、对研究结论边界的审慎，都让我深刻认识到学术训练不仅在于完成一项工作，更在于培养一种对问题负责、对事实负责、对自己负责的态度。能够在硕士阶段师从张老师，是我求学道路上的幸运，也是我十分珍惜的一段经历。张老师给予我的影响，早已超出这篇论文本身，并将长久地影响我今后的学习、工作与人生。

同时，我也衷心感谢课题组各位老师和实验室同学们在学习与科研过程中给予我的帮助。感谢大家在文献阅读、技术讨论和实验推进中带来的启发与支持，让我在不断交流中逐步完善自己的思考。感谢我的家人始终如一的理解、陪伴与守护，是你们的信任和支持给了我安心前行的力量。

最后，再次向所有关心、帮助和陪伴过我的人致以诚挚谢意。谨以此文，献给在我成长道路上给予我温暖与力量的每一个人。

\newpage
\fi
